\begin{circuitikz}[scale=0.92, every node/.style={transform shape}, font=\sffamily\small]
% ================= POW-BB rail splitter
\draw[rounded corners=2pt, fill=hwcol] (0,2.4) rectangle (3.2,5.8);
\node[font=\sffamily\bfseries\small] at (1.6,5.45) {SBC-POW-BB};
\node[font=\sffamily\scriptsize] at (1.6,5.1) {labelled rail splitter};
\foreach \y/\name in {4.4/3V3 rail, 3.7/5V rail, 3.0/GND}{
  \draw (3.2,\y) -- (3.8,\y);
  \node[anchor=east, font=\sffamily\scriptsize] at (3.1,\y) {\name};
}
% ================= MCC 118 with its terminal block
\draw[rounded corners=2pt, fill=hwcol] (7.0,1.6) rectangle (11.6,5.8);
\node[font=\sffamily\bfseries\small] at (9.3,5.45) {MCC 118, board 0};
\node[font=\sffamily\scriptsize] at (9.3,5.1) {8 ch, 12-bit, 100 kS/s};
\node[font=\sffamily\scriptsize] at (9.3,4.75) {address jumpers all OFF};
% screw terminals drawn as small squares on the left edge
\foreach \y/\name in {4.4/CH0, 3.7/CH1, 3.0/AGND, 2.2/DGND}{
  \draw[fill=gray!35] (6.8,\y-0.13) rectangle (7.0,\y+0.13);
  \node[anchor=west, font=\sffamily\scriptsize] at (7.1,\y) {\name};
}
\node[anchor=west, font=\sffamily\scriptsize] at (7.1,1.9) {CH2 to CH7 unused};
% the three rail wires, straight and level
\draw (3.8,4.4) -- (6.8,4.4);
\draw (3.8,3.7) -- (6.8,3.7);
\draw (3.8,3.0) -- (6.8,3.0);
\node[font=\sffamily\scriptsize, fill=white, inner sep=1.5pt] at (5.3,4.58) {expect about 3.30 V};
\node[font=\sffamily\scriptsize, fill=white, inner sep=1.5pt] at (5.3,3.88) {5 V, only because $5 < 10.1$};
\node[font=\sffamily\scriptsize, fill=white, inner sep=1.5pt] at (5.3,3.18) {common analog ground};
% ================= the Pi under the HAT
\draw[rounded corners=2pt, fill=extcol] (7.0,-2.3) rectangle (11.6,0.4);
\node[font=\sffamily\bfseries\small] at (9.3,0.05) {Raspberry Pi 4};
\node[font=\sffamily\scriptsize, align=left, anchor=north west] at (7.15,-0.3)
  {SPI0 CE0 on pin 24\\SCLK 23, MOSI 19, MISO 21\\address GPIO12 / 13 / 26\\on pins 32 / 33 / 37\\ID EEPROM on the HAT ID pins};
\draw[very thick] (9.3,0.4) -- (9.3,1.6);
\node[anchor=west, font=\sffamily\scriptsize] at (9.45,1.0) {40-pin header, this HAT alone};
% DGND back to the Pi ground
\draw (6.8,2.2) -- (5.6,2.2) -- (5.6,-0.9) -- (7.0,-0.9);
\node[font=\sffamily\scriptsize, anchor=west] at (5.75,1.35) {DGND to Pi GND};
\node[ground] at (5.6,-0.9) {};
% ================= the limit
\draw[rounded corners=2pt, fill=red!5, draw=red!55!black] (0,-1.6) rectangle (4.6,1.0);
\node[font=\sffamily\bfseries\small, anchor=west, text=red!60!black] at (0.15,0.65) {Input limits};
\node[font=\sffamily\scriptsize, align=left, anchor=north west, text=red!60!black] at (0.15,0.3)
  {Never exceed +/-10.1 V on any input.\\Pi GPIO is not an analog source.\\Warm up 1 minute before quoting\\a number.};
\end{circuitikz}
